\section{The angular clearance source as a weighted unstable family}
\label{sec:ag-source}
The proofs in Sections~\ref{sec:xf-jacobi}--\ref{sec:xf-pressure}
are inherited, including their stated pressure premise. We now give a
different summation route. It starts from the actual angular source
in the one-mark flux, rather than taking a supremum over every
physical cylinder. The geometric results of this section use no
pressure or growth assumption. The propagation theorem in the next
section uses the explicitly specified homogeneous curve-growth input,
not the stronger physical-word pressure input of
Input~\ref{input:v112-geometric-pressure}.

All radii, sections, time conventions and exact records are unchanged.
Put $b_*=R_++\delta_*=97/200$, $a_*=a_\circ=\sqrt{1827}/200$,
$\kappa_-=1/R_+$ and $\kappa_+=1/R_-$. We use actual boundary
arclength $s=R\alpha$ and outgoing angle $\varphi$, with
$c=\cos\varphi>0$ and $p=\sin\varphi$. Arclength on a collision
curve always means Euclidean arclength in $(s,\varphi)$.

\subsection{A genuine first collision, not a reflection at the mark}
Fix $0<u\le\delta_*$ and $\sigma\in\{-1,1\}$, and put
$b=\sigma(R+u)$. For $v_\vartheta=(\cos\vartheta,\sin\vartheta)$,
the marked foot is $q_f=-b v_\vartheta^\perp$. Starting at this
point with velocity $v_\vartheta$, let $q_1=q_f+a v_\vartheta$
be the first actual collision, and let $Y_{R,\sigma,u}(\vartheta)$
be the outgoing state there. Thus, in the one-mark notation,
\begin{equation}\label{eq:ag-first-image}
 Y_{R,\sigma,u}(\vartheta)=T_RX_R(\vartheta,\sigma,u).
\end{equation}
The line is exterior to the marked disk. Moreover, for every other
lattice center $z$, $|z-q_f|\ge1-b_*>R_+$. The foot is therefore
in the actual free region for every $\vartheta$. The inherited
finite-horizon bound supplies the two incident collisions. We exclude
the finitely many first-collision tangencies when differentiating.
The marked disk itself is never reflected from at the foot.

\begin{theorem}[Exact first-image density and uniform geometric regularity]
\label{thm:ag-density}
On every regular interval of fixed first collision center, let $s$
be an unwrapped boundary coordinate at that center. Then
\begin{align}
 \frac{ds}{d\vartheta}&=-\frac{a}{c},&
 V:=\frac{d\varphi}{ds}&=\frac1R+\frac{c}{a},&
 \rho:=\left|\frac{d\vartheta}{d\ell}\right|
     &=\frac{c}{a\sqrt{1+V^2}}.\label{eq:ag-density}
\end{align}
In particular, with
\begin{align}
 U_0&=\kappa_++a_*^{-1},&
 K_0&=U_0/a_*+(1+b_*/a_*)/a_*^2,\label{eq:ag-geometry-constants}
\end{align}
every such image is an unstable graph satisfying
\begin{equation}\label{eq:ag-regularity}
 a\ge a_*,\qquad \kappa_-\le V\le U_0
 <\kappa_++\tau_-^{-1},\qquad
 |dV/ds|\le K_0,\qquad 0<\rho\le c/a_*.
\end{equation}
These bounds hold at every regular first collision, with no lower
cutoff on $c$ and no lower cutoff on $u$ within its positive range.
\end{theorem}
\begin{proof}
The foot standoff is Lemma~\ref{lem:xf-standoff}. Let $n$ be the
outward normal at $q_1$ and $t=n^\perp$. The incoming and outgoing
velocities are $v^-=-cn+pt$ and $v^+=cn+pt$. At fixed $u$, the
variation of $q_f$ is parallel to $v^-$. Differentiating
$q_1=q_f+a v^-$ and projecting onto $(v^-)^\perp$ gives
$-c\,ds/d\vartheta=a$, proving the first formula. The specular
angle identity gives
$d\vartheta^+=2R^{-1}ds-d\vartheta$; subtracting the normal-angle
differential $R^{-1}ds$ yields
$d\varphi=R^{-1}ds-d\vartheta$. Division by $ds$ proves the
formula for $V$. Since $d\ell=\sqrt{1+V^2}|ds|$, the density
formula follows. Its vanishing factor $c$ is essential.

For clarity, the second-derivative estimate does not appeal to a
sampled trajectory or to bounded curvature by name. Projecting the
same contact identity onto $n$ gives
$da/d\vartheta=-b-ap/c$. Hence
\begin{equation}\label{eq:ag-a-derivative}
 \frac{da}{ds}=p+\frac{bc}{a},\qquad
 \frac{dc}{ds}=-pV,\qquad
 \frac{dV}{ds}=-\frac{pV}{a}
               -\frac{c}{a^2}\left(p+\frac{bc}{a}\right).
\end{equation}
Use $|p|,c\le1$, $|b|\le b_*$ and $a\ge a_*$ to obtain
\eqref{eq:ag-regularity}. The curvature of the graph in
$(s,\varphi)$ is $|dV/ds|/(1+V^2)^{3/2}$ and is therefore also
bounded by $K_0$. Finally $a_*>\tau_-$ gives the strict upper-cone
inequality. No finite derivative at $c=0$ is asserted.
\end{proof}

\subsection{Uniform finite geometry before the homogeneous cut}
Use the finite candidate set already fixed in
\eqref{eq:xf-candidates}, and write $N_*=N_\circ$ in this section.
For reference, its radius is $D_*=\tau_++2R_++\delta_*$. It has the
explicit packing upper bound
\begin{equation}\label{eq:ag-packing}
 N_*+1\le4(D_*+1/2)^2.
\end{equation}
Indeed, the disjoint open radius-$1/2$ disks centered at all the
lattice points in this closed radius-$D_*$ disk lie in the
radius-$(D_*+1/2)$ disk. This bound does not assign a sampled value
to $\tau_+$. Here $\tau_+$ means the inherited uniform supremum of
actual free-flight lengths over the original radius family.

\begin{lemma}[A uniform branch count]
\label{lem:ag-branches}
For each fixed $R,u,\sigma$, the angular circle can be cut at at
most $8N_*$ points, in addition to an angular coordinate seam, so
that on every resulting interval the first collision center is
fixed, the contact coordinate has no seam, and the image in
Theorem~\ref{thm:ag-density} is injective. The number of intervals
is at most $M_*=8N_*+1$; each has image length at most
$2\pi R_+\sqrt{1+U_0^2}$.
\end{lemma}
\begin{proof}
Every first collision center belongs to the stated candidate set:
its distance from the marked center is at most
$\tau_++2R_++\delta_*$. For a candidate $z$, loss or creation of
an intersection of the line with its disk requires
\[
                 z\cdot v_\vartheta^\perp+b=R
        \quad\hbox{or}\quad z\cdot v_\vartheta^\perp+b=-R.
\]
Each equation has at most two roots on the angular circle, since
$z\ne0$. There are thus at most $4N_*$ tangency cuts. Away from
these cuts, positive first-hit times vary continuously. Their
ordering cannot change without two distinct disks having a common
contact point, which is impossible because $2R<1$. No intersection
can pass through the foot, which is exterior to every disk as
proved above. Thus the first target is fixed on each interval.

At each candidate disk choose one contact-coordinate seam, at
$q=z+R(1,0)$. A line in the source passes through this point only
if $q\cdot v_\vartheta^\perp=-b$. Since
$|q|\ge1-R_+>b_*$, there are at most two such roots for each disk.
We allow $4N_*$ seam cuts in the displayed bound, which is a
harmless overestimate. Coincident roots only reduce the count.
On a seam-free interval $ds/d\vartheta=-a/c<0$, so $s$ is strictly
monotone and cannot wind around its contact circle. Its range has
length at most $2\pi R$. This proves injectivity and the length
bound. The argument uses no long collision-word count. Overlapping
images from different intervals remain separately labeled sources;
they are not silently identified in the quotient collision space.
\end{proof}

\subsection{Paying the infinitely many grazing pieces}
Fix $k_0\ge2$. Define the homogeneous strips
\[
 H_k^+=\{\pi/2-k^{-2}<\varphi<\pi/2-(k+1)^{-2}\},\quad
 H_k^-=-H_k^+,\qquad k\ge k_0,
\]
and let $H_0$ be the remaining central strip, apart from its boundary.
Let $\ell_*>0$ be a fixed allowed curve length. Subdivide each
homogeneous image into pieces of length at most $\ell_*$. Set
\begin{align}
 L_*&=1+\left\lceil\frac{2\pi R_+\sqrt{1+U_0^2}}{\ell_*}\right\rceil,
 & A_*&=\frac{2M_*L_*}{a_*}
                    \left(1+\frac{2}{k_0-1}\right).
                                                  \label{eq:ag-budget}
\end{align}
The factor two in $A_*$ includes both first-impact signs.

\begin{theorem}[A uniformly summable angular source]
\label{thm:ag-family}
For every $R\in[R_-,R_+]$ and $0<u\le\delta_*$, the pushforward
of $\sum_{\sigma=\pm1}d\vartheta$ under the actual maps
$Y_{R,\sigma,u}$ has a countable, source-labeled representation
\begin{equation}\label{eq:ag-family}
       \mu_{R,u}=\sum_{C\in\mathcal F_{R,u}}\rho_C\,d\ell_C,
       \qquad \sum_{C\in\mathcal F_{R,u}}\|\rho_C\|_\infty\le A_*.
\end{equation}
Each $C$ is a homogeneous unstable graph of length at most $\ell_*$,
with the slope and curvature bounds in \eqref{eq:ag-regularity}.
The identity holds against every nonnegative Borel test, and has
total mass $4\pi$. An arbitrary Borel selector in the original
angular source can be retained on the right as a pullback selector;
its derivatives or connected components are not counted.
\end{theorem}
\begin{proof}
On each interval of Lemma~\ref{lem:ag-branches}, $V>0$ and $s$
is strictly monotone, so $\varphi$ is strictly monotone. Each
homogeneity strip is consequently visited at most once by that
interval. The length subdivision requires at most $L_*$ pieces
per strip. The density is exactly \eqref{eq:ag-density}, not a
constant-density substitute. On $H_0$ it is at most $1/a_*$. On
either $H_k^\pm$, it is at most $k^{-2}/a_*$, because
$\cos\varphi\le k^{-2}$. Thus, summing over both signs,
\[
 \sum_C\|\rho_C\|_\infty
 \le\frac{2M_*L_*}{a_*}
                   \left(1+2\sum_{k=k_0}^\infty k^{-2}\right)
 \le A_*.
\]
The last inequality follows by integrating $x^{-2}$ from $k_0-1$
to infinity. First perform the exact substitution on finitely many
regular intervals and truncated homogeneous strips. The omitted
angular cut points are finite, and strip boundaries have at most
one preimage per regular interval. Monotone exhaustion proves the
measure identity and its total mass. Borel selectors are multiplied
before substitution and are not differentiated. In particular an
infinite number of small selector components does not incur an
infinite multiplicity charge.
\end{proof}

Theorem~\ref{thm:ag-family} is the new first-source result. Replacing
$\rho_C$ by the crude bound $1/a_*$ \emph{before} summing the
homogeneity strips would give a divergent series. The factor $c$
is the geometric payment for that series. This is a source estimate
in a specific curve norm, not a statement that $\mu_{R,u}$ has a
two-dimensional collision density or a bounded original roof density.

